\Askhsh{22046}{2}
\begin{ylh}{1}
    \item 9.2. Το Πυθαγόρειο θεώρημα
    \item 11.5. Μήκος τόξου.
\end{ylh}
\begin{wrapfigure}[11]{r}{5.5cm}
\centering
\begin{tikzpicture}[scale=1.5]
    % Ορισμός σημείων
    \tkzDefPoint(0,0){O}
    \tkzDefPoint(1,0){Gamma} % Γ
    \tkzDefPoint(2,0){B}
    % Σημείο A: Τρίγωνο OAB είναι ορθογώνιο στο A. OA=R=1, OB=2R=2. Άρα η γωνία OBA=30°.
    % Η γωνία BOA είναι 60°. Θέτουμε το Α στις συντεταγμένες (R*cos(60°), R*sin(60°))
    \tkzDefPoint(0.5,{sqrt(3)/2}){A}

    % Σχεδίαση κύκλου
    \tkzDrawCircle[R](O,1 cm) % Κύκλος με κέντρο O και ακτίνα 1cm

    % Σχεδίαση τμημάτων
    \tkzDrawSegments(O,A A,B B,Gamma)
    \tkzDrawSegment[thick](O,B) % Το OB είναι η προέκταση του OΓ

    % Σημάνσεις
    \tkzMarkSegments[mark=s|,size=3pt](O,Gamma Gamma,B O,A) % Σημειώνουμε OΓ, ΓB, OA ως ίσες (R=1)
    \tkzLabelSegment[left](O,A){R=1} % Ετικέτα OA με R=1

    % Ετικέτες σημείων
    \tkzLabelPoint[left](O){O}
    \tkzLabelPoint[right](B){B}
    \tkzLabelPoint[below left](Gamma){$\Gamma$}
    \tkzLabelPoint[above right](A){A}

    % Σήμανση ορθής γωνίας στο A
    \tkzMarkRightAngle(O,A,B)

    % Σκιασμένη περιοχή: τμήματα AB, BΓ και τόξο AΓ
    \fill[maincolor!30] (B) -- (A) arc (60:0:1cm) -- cycle;
\end{tikzpicture}
\end{wrapfigure}
Δίνεται κύκλος με κέντρο Ο και ακτίνα $R = 1$. Θεωρούμε ακτίνα Ο$\Gamma$ την οποία προεκτείνουμε κατά τμήμα $\Gamma \text{Β} = \text{Ο}\Gamma = R$ και το εφαπτόμενο τμήμα $\text{ΒΑ}$, όπως φαίνεται στο σχήμα.

\begin{alist}
\item Να αποδείξετε ότι $\text{Ο}\widehat{\text{Β}}\text{Α} = 30^{\circ}$. \monades{5}
\item Να αποδείξετε ότι $\text{ΑΒ} = \sqrt{3}$. \monades{10}
\item Να υπολογίσετε την περίμετρο του γραμμοσκιασμένου μικτόγραμμου τριγώνου $\text{ΑΒ}\Gamma$. \monades{10}
\end{alist}